% Plan: development/outline.md, section 6, item B. The family files B1-B9
% and B10-split-extras are subsections of this section (Sections S1.1 to
% S1.10 of the supplement).
\section{Certificates by family}\label{app:families}

Family by family, we complete the results that \cref{sec:split,sec:other} state with a proof sketch or without proof: the details of the stored model that the proofs use, the full proofs, and the computations with their inputs and numbers.
Each result has one certificate box, with six fields: inputs (SHA-256 prefixes), computation and arithmetic, data reading, implementations, evidence level and replay tier (\cref{sec:semantics-trust,sec:semantics-protocol,sec:repro}); in this supplement, the box is the one place that states how a result was verified.
Where the main text already proves a result in full, as for \cref{thm:lnts-opt,prop:dtoc5-identity,thm:dtoc5-bracket}, the subsection adds only the data and the computations.
The listed status and the prior results are in \cref{tab:closures,sec:results-prior,app:literature}; we repeat them only where an argument uses them.
All statements concern \reading{b} (\cref{sec:semantics-model}).
\emph{The first code} of a certificate is its first implementation and \emph{the second code} its second implementation; \cref{sec:semantics-protocol} states that AI agents wrote both, what ``separately written'' means and the resulting common-mode risk.
The proofs of the validity results of \cref{sec:split-framework} are in \cref{app:splitproofs}, two explanatory results on splits that no certificate uses in \cref{app:split-explanatory}, the exactly feasible points in \cref{app:points}, and the rounding-error analysis of the \inst{eg} certificates in \cref{app:egrounding}.
Run records and superseded results are in the artifact (\cref{app:repro-register}).
\Cref{tab:staged} summarizes the fifteen staged certificates of \cref{sec:split} and the \inst{waterno2_06} bound.

% Generated by paper-open-minlplib/data/make_tables.py from data/numbers.json.
% Do not edit by hand; rerun the script. Requires booktabs, array and rotating;
% \inst{} typesets an instance name; underscores are not escaped (macros.tex).
\begin{sidewaystable}
\centering
\footnotesize
\setlength{\tabcolsep}{3pt}
\caption{The 15 staged certificates and the \inst{waterno2_06} bound.
Stages and separator dimension; free variables (no finite bounds in the model); split form; validity result; window or case split;
how each stage problem is minimized and the branching dimension; arithmetic tags of the dual (\cref{sec:semantics-trust}).}
\label{tab:staged}
\begin{tabular}{@{}>{\raggedright\arraybackslash}p{2.4cm}>{\raggedright\arraybackslash}p{2.6cm}>{\raggedright\arraybackslash}p{2.2cm}>{\raggedright\arraybackslash}p{3.0cm}>{\raggedright\arraybackslash}p{2.4cm}>{\raggedright\arraybackslash}p{2.6cm}>{\raggedright\arraybackslash}p{3.2cm}>{\raggedright\arraybackslash}p{1.0cm}@{}}
\toprule
instances & stages; separator & free variables & split form & validity & window or case split & stage check; branching & arith. \\
\midrule
\inst{lnts50}--\inst{lnts400} & $N$ steps; states eliminated & states except the fixed end values & affine, for each fixed $h$ & \cref{lem:split-bound}(d) & global $h$ by monotonicity & closed form; none & \arith{E} \\
\inst{dtoc5} & $T=49{,}999$; dim.\ 1 & all except $y_0$ & affine (costates) & \cref{lem:split-bound}(a) & none & closed-form convex quadratics; none & \arith{E} \\
\inst{optcdeg2} & $N=50{,}000$; dim.\ 2 & $y_1,\dots,y_N$ & quadratic in $v$, $q_t=0$ at both switches & \cref{lem:split-bound}(a) with state enclosures & none & exact (Sturm); 1-D cells in $v$ & \arith{E} \\
\inst{lukvle10} & 497 pairs + tail; dim.\ 2 & all & affine (pair Lagrangian) & \cref{lem:split-bound}(d); row form of \cref{prop:split-window} & 3-pair tail, 2-D entry state & interval B\&B; 2-D & \arith{I} \\
\inst{chain50}--\inst{chain400} & $N$ stages; lifted dim.\ 2 & all except ends & field type (discrete catenary) & \cref{lem:split-bound}(c), (a) & end values (2-D) & interval B\&B; 2-D & \arith{I} \\
\inst{catmix100}--\inst{catmix800} & $N$ stages; ray dim.\ 1 & states except $x_0$ & concave chord minorants & \cref{lem:split-minorant} & none & per-ray interval bounds; 1-D & \arith{F} \\
\inst{waterno2_06} (not closed) & 6 periods of 166 variables; dim.\ 3 & monomial auxiliaries only & cellwise affine & \cref{prop:split-cellwise} & none & B\&B per cell pair & \arith{F}/\arith{E} \\
\bottomrule
\end{tabular}
\end{sidewaystable}

